\begin{recipe}{Coleslaw}
\kicker[Bijgerecht,Vegetarisch,Amerikaans]{Bijgerecht · vegetarisch · Amerikaans}
\index[register]{Coleslaw}
\index[register]{Witte kool}
\index[register]{Rodekool}
\index[register]{Wortel}
\index[register]{Mayonaise}

\meta{15 minuten + 1 uur rusten \dvd Voor 4 tot 6 personen \dvd Eenvoudig}

\lettrine{D}{eze} Amerikaanse coleslaw is een frisse, romige koolsalade van rauwe
witte en rode kool met wortel. Hij is snel gemaakt, maar wordt lekkerder als je
hem minstens een uur laat intrekken in de koelkast. Een vaste metgezel bij
barbecue of pulled pork.

\blockrule{Ingrediënten}
\begin{ingredients}
  \ing{300 g}{witte kool (of spitskool), flinterdun gesneden}\index[register]{Witte kool}
  \ing{150 g}{rode kool, flinterdun gesneden}\index[register]{Rodekool}
  \ing{1}{grote winterpeen, geraspt of in dunne reepjes}\index[register]{Wortel}
  \ing{1}{kleine rode ui, heel fijn gesneden (optioneel)}\index[register]{Rode ui}
  \ing{5 el}{mayonaise}\index[register]{Mayonaise}
  \ing{2 el}{volle yoghurt of zure room}
  \ing{2 el}{natuurazijn of appelciderazijn}
  \ing{1 el}{kristalsuiker of honing}
  \ingb{peper en zout, naar smaak}
\end{ingredients}

\blockrule{Bereiding}
\tip{Wil je dat de kool extra knapperig blijft? Bestrooi de gesneden kool met een
theelepel zout, laat het 30 minuten in een vergiet uitlekken en knijp het vocht
eruit voordat je de dressing toevoegt. Afgedekt is de coleslaw zo'n 2 tot 3
dagen houdbaar in de koelkast; roer hem voor het serveren nog even door.}
\begin{steps}
  \item Was de witte kool, rode kool en winterpeen. Snijd de witte kool en rode
        kool flinterdun en rasp de winterpeen, of snijd hem in dunne reepjes.
        Meng in een grote kom met de rode ui.
  \item Klop in een aparte kom de mayonaise, yoghurt (of zure room), azijn en
        suiker glad. Breng op smaak met peper en zout.
  \item Schep de dressing door het koolmengsel tot alles goed bedekt is.
  \item Dek de kom af en zet de coleslaw minstens 1 uur in de koelkast, zodat
        de kool wat zachter wordt en de smaken intrekken.
\end{steps}
\end{recipe}
